
\[
\alpha:P\xrightarrow{i}A\xrightarrow{j}A\otimes_{\mathbb Z}\mathbb Q
\]
\uline{(où} $j(x)=x\otimes1$\uline{) est une application polynôme. Si} $I^{m+1}=0$\uline{, alors} $\alpha$ \uline{est de degré} $\leq m$.

\emph{Démonstration} : Soit $x_1,\ldots,x_r$ une suite d'éléments de $F$ avec $x_i=1+y_i$ ($i=1,\ldots,r$). Alors la formule binomiale donne pour tout $n_1,\ldots,n_r\in Z$ :
\begin{equation}\tag{2.1.5}
x_1^{n_1}\cdots x_r^{n_r}=
\sum_{k_1,\ldots,k_r\geq0}\binom{n_1}{k_1}\cdots\binom{n_r}{k_r}y_1^{k_1}\cdots y_r^{k_r},
\end{equation}
où la somme à droite est finie puisque les éléments $y_1,\ldots,y_r$ sont nilpotents. Mais les coefficients $\binom{n_1}{k_1}\cdots\binom{n_r}{k_r}$ sont des polynômes à coefficients rationnels dans les variables $n_1,\ldots,n_r$ de degré $k_1+\cdots+k_r$, de sorte que $\alpha$ est bien une application polynôme. Si $I^{m+1}=0$, alors $y_1^{k_1}\cdots y_r^{k_r}=0$ pour $k_1+\cdots+k_r>m$. Donc $\alpha$ est de degré $\leq m$, C.Q.F.D.

\subsection*{2.2.}
Considérons maintenant un schéma noethérien $X$ et un élément $y$ de $\Filt_j(X)$ (1.1.1). En associant à tout $x\in K^{\bullet}(X)$ l'élément $xy$ de $K_\bullet(X)$ on obtient un homomorphisme de $K^{\bullet}(X)$-modules
\[
\Theta_y:K^{\bullet}(X)\longrightarrow K_\bullet(X)
\]
qui se factorise, compte tenu de (1.3.2), par l'épimorphisme canonique
\[
j:K^{\bullet}(X)\longrightarrow K^{\bullet}(X)/\Filt^{j+1}(X)=A.
\]
